\section{Purpose, source audit, and precise claims}
\label{sec:scope}

This report develops a quantitative interface for continuous mathematics in
intuitionistic primitive recursive arithmetic (\ipra). Its main new work,
relative to the inspected ProveIt packets, concerns integration and
transport: how to construct an integrable limit from weaker observations;
how to reconstruct a density from coherent finite mass data; how to carry
errors through operators, averaging, and topological approximation; and
which substitutions preserve the chosen observations. These questions are
connected. An operation is usable only after the observation in which it is
stable has been specified.

The source baseline is ProveIt commit
\texttt{b2341af5370cd8a43f72ebb2341c0d7d7395b634}, inspected on
20 November 2026. The source tree was obtained through GitHub, and text
files were cross-checked against an existing local checkout at that exact
commit. No remote repository mutation is part of this delivery.
The repository itself separates the kernel-checked libraries
\texttt{ConciseIPRA}, \texttt{IPRADerivations}, and \texttt{IPRABridge}
from delivered research packets that have not necessarily been verified
\cite{repo-main,repo-unification}.

\subsection{What is inherited, what is added}

The November 20 continuous-framework packets already develop regular
names, completion, uniform extension, finite-net semantics, finite
gluing, and several analytic constructions
\cite{repo-continuous,repo-continuous-two}. The probability packet already
contains finite expectation, effective integration, conditional expectation,
and convergence statements with quantitative assumptions
\cite{repo-expectation}. We do not rename these constructions and present
them as new discoveries. The present contribution is a collection of
sharper interfaces and explicit constructions that connect those areas.

\begin{enumerate}
\item A corrected, event-free truncation inequality upgrades a Cauchy
sequence for the bounded metric
\(\rho(f,g)=\int\min(1,|f-g|)\) to a Cauchy sequence in \(L^1\), given a
uniform clipping-tail modulus. It constructs an actual regular name;
the limit need not be supplied in advance.
\item Coherent dyadic masses produce a density from an energy-tail
modulus. A bound on sibling oscillation supplies that modulus explicitly.
An energy bound by itself gives a primitive-recursive finite stability
window, with a different conclusion from a limit name.
\item Completed bounded operators, product integration, and expectation
quantifiers share the same norm estimates. Arbitrary first-order
substitution is not sound for an \(L^1\) semantics; bounded-compression
pullbacks and an explicit Lipschitz upgrade repair specific cases.
\item Rational distance intervals produce interleavings of Rips
filtrations and finite homological certificates without deciding unstable
distance-threshold equalities.
\item Integral equations with a certified contraction gap and analytic
transforms of rough \(L^1\) data are built by the same calculus. Backward
error propagation permits reuse of one approximation across many
observations.
\end{enumerate}

The classical ideas behind completion, Vitali convergence, conditional
expectation, and persistence stability are established mathematics.
Constructive integration through completion has a substantial literature
\cite{bishopcheng,petrakiszeuner}. Our claims are explicit constructions,
corrected proof steps, finite estimates, and an \ipra\ implementation
discipline. We make no priority claim for the classical theorems.

\subsection{A concrete repair in the inspected source}

The separate-tail shortcut displayed in the proof of the ``Vitali form''
in the probability packet's Section 12 is false. Its proposed bound has
the form
\[
 |x-y|\leq |x-y|\mathbf1_{\{|x|\leq K,\ |y|\leq K\}}
       +|x|\mathbf1_{\{|x|>K\}}+|y|\mathbf1_{\{|y|>K\}}.
\]
For \(x=2\), \(y=-1\), and \(K=1\), its sides are \(3\) and \(2\).
The desired convergence theorem is repairable: clip both variables
before comparing them, and retain both clipping residuals. The repaired
proof appears below, with the displayed counterexample in the exact
arithmetic checks.

This defect does not invalidate the library's different, correctly stated
\emph{common-envelope} truncation theorem. In particular,
\texttt{Probability.TruncationEstimate} already contains object-\ipra
derivations for a bound under \(|X-Y|\leq2G\). The source audit and Lean
roadmap distinguish that existing result from the separate-tail statement
repaired here.

\subsection{Four kinds of evidence}

Throughout the report we distinguish:
\begin{description}
\item[Human finite proofs.] Complete mathematical arguments whose
operations, bounds, and induction invariants are given explicitly.
\item[Checked host-Lean lemmas.] The accompanying small Lean modules
prove rational clipping and energy-search lemmas in Lean's metatheory.
They have been compiled. Their types are not \texttt{Prf PRAAx}.
\item[Object-theory targets.] A Lean definition of an arithmetic sentence
and a plan to construct \(\Gamma\vdash[\mathrm{PRAAx}]\varphi\).
The new analytic endpoints have not all been reified as such derivations.
\item[External interpretation.] In ordinary mathematics a regular name
denotes an element of a completion. This explains its familiar meaning;
it is not an axiom of \ipra.
\end{description}

An exact-arithmetic experiment is useful for exposing bad constants or
indexing, but does not replace any of these proofs. The validation ledger
states exactly what was run. In particular, this report does not claim a
uniform embedding of unrestricted classical analysis, arbitrary measure
spaces, classical compactness, or unrestricted higher-order logic into
\ipra.
